\documentclass{article}

\usepackage[a4paper, left=1.5cm, right=1.5cm, top=2cm, bottom=2cm]{geometry}

\usepackage{../../../../components/components}

\usepackage{fancyhdr}


% Configuration des en-têtes et pieds de page
\pagestyle{fancy}
\fancyhf{} % reset tout

\fancyhead[L]{DL2 Math-Info PR4}
\fancyhead[C]{Probabilités}
\fancyhead[R]{2025-2026}

\fancyfoot[L]{Ewen Rodrigues de Oliveira}
\fancyfoot[R]{\thepage}

\begin{document}

\docTitle{Chapitre 6 : Vecteurs aléatoires}

\section{Covariance}

\definition{
    Soient $X,Y$ deux variables aléatoires réelles discrètes de carrés intégrables. On appelle \textbf{covariance} de $X$ et $Y$ la quantité suivante :
    \[\text{Cov}(X,Y) = \mathbb{E}[(X-\mathbb{E}[X])(Y-\mathbb{E}[Y])] = \mathbb{E}[XY] - \mathbb{E}[X]\mathbb{E}[Y]\]
}

\remark{Souvent, si deux variables ont tendance à être plus grandes que leurs moyennes en même temps (respectivement plus petites), on aura $Cov(X,Y) > 0$. Si au contraire, lorsque l'une est plus grande que sa moyenne, l'autre a tendance à être plus petite, on aura $Cov(X,Y) < 0$.} 

\example{
    On tire un dé standard avec des faces numérotées de 1 à 6. On pose :
    \begin{itemize}
        \item $X = 1$ si le résultat est pair et 0 sinon,
        \item $Y = 1$ si on a tiré un 6 et 0 sinon,
        \item $Z = 1$ si on a tiré un 1 et 0 sinon.
    \end{itemize}
    Ces trois variables sont toutes différentes. Après calcul, on trouve que $\mathbb{E}[X] = \frac{1}{2}$, $\mathbb{E}[Y] = \frac{1}{6}$, $\mathbb{E}[Z] = \frac{1}{6}$.\\
    De plus, on remarque que $XY = Y$ (car $Y=1$ implique $X=1$), $XZ = 0$ et $YZ =0$ (car les variables ne peuvent pas valoir $1$ en même temps).\\
    On a donc : 
    \[
    \text{Cov}(X,Y) = \mathbb{E}[XY] - \mathbb{E}[X]\mathbb{E}[Y] = \frac{1}{6} - \frac{1}{2}\cdot\frac{1}{6} = \frac{1}{12}
    \]
    Ainsi que :
    \[
    \text{Cov}(X,Z) = -\frac{1}{12} \quad \text{et} \quad \text{Cov}(Y,Z) = -\frac{1}{36}
    \]
    Ainsi, $X$ et $Y$ sont positivement corrélés, tandis que $X$ et $Z$ ainsi que $Y$ et $Z$ sont négativement corrélés.
}

\problem{On retrouve le même problème que pour la variance quand on s'intéresse aux unités. Par exemple, si $X$ et $Y$ sont des variables aléatoires mesurant des longueurs en mètres, alors $\text{Cov}(X,Y)$ aura des unités de mètres carrés, ce qui peut rendre l'interprétation difficile.}

\definition{
    Soient $X,Y$ deux variables aléatoires réelles discrètes de carrés intégrables. On appelle \textbf{coefficient de corrélation} de $X$ et $Y$ la quantité suivante :
    \[\rho(X,Y) = \frac{\text{Cov}(X,Y)}{\sqrt{\text{Var}(X)}\sqrt{\text{Var}(Y)}}\]
}

\remark{Le coefficient de corrélation $\rho(X,Y)$ est un nombre sans unité qui mesure la force et la direction de la relation linéaire entre $X$ et $Y$. Il est compris entre -1 et 1, où 1 indique une corrélation positive parfaite, -1 une corrélation négative parfaite, et 0 aucune corrélation linéaire.}

\theorem{Proposition}{Identités de la variance et covariance}{false}{
    Soient $X, X', Y, Y'$ des variables aléatoires réelles de carrés intégrables et $a, b \in \mathbb{R}$. On a :
    \begin{enumerate}
        \item $\text{Cov}(X,X) = \text{Var}(X)$
        \item $\text{Cov}(X+a,Y+b) = \text{Cov}(X,Y)$
        \item $\text{Cov}(aX,bY) = ab\text{Cov}(X,Y)$
        \item $\text{Var}(aX+b) = a^2\text{Var}(X)$ (ne pas confondre avec $\mathbb{E}[aX+b] = a\mathbb{E}[X]+b$)
        \item $\text{Cov}(X,Y) = \text{Cov}(Y,X)$
        \item $\text{Cov}(X+X',Y+Y') = \text{Cov}(X,Y) + \text{Cov}(X,Y') + \text{Cov}(X',Y) + \text{Cov}(X',Y')$
        \item $\text{Var}(X+X') = \text{Var}(X) + 2\text{Cov}(X,X') + \text{Var}(X')$
    \end{enumerate}
}

\training{Démontrer les différentes propriétés de la proposition précédente.}
\noindent \carreaux{10}

\theorem{Proposition}{Lien entre covariance et indépendance}{false}{
    Soient $X,Y$ deux variables aléatoires réelles de carrés intégrables. On a :
    \[
        X \text{ et } Y \text{ sont indépendantes} \Rightarrow \text{Cov}(X,Y) = 0
    \]
}

\attention{La réciproque n'est pas vraie : il est possible que $\text{Cov}(X,Y) = 0$ alors que $X$ et $Y$ ne sont pas indépendantes.}

\cexample{
    Soit $X$ une variable aléatoire de loi de Bernoulli de paramètre $\frac{1}{2}$, et posons $Y = X^2$. Alors $X$ et $Y$ ne sont pas indépendantes, mais $\text{Cov}(X,Y) = 0$.
}

\theorem{Corollaire}{Variance d'une somme de variables aléatoires}{false}{
    Soient $X_1, \ldots, X_n$ des variables aléatoires réelles de carrés intégrables. On a :
    \[
        Var(\sum_{i=1}^n X_i) = \sum_{i=1}^n Var(X_i) + 2\sum_{1 \leq i < j \leq n} Cov(X_i, X_j)
    \]
    De plus, si les variables sont indépendantes, alors : $Var(\sum_{i=1}^n X_i) = \sum_{i=1}^n Var(X_i)$.
}

\example{
    Soit $X_1, \ldots, X_n$ des variables aléatoires indépendantes de loi de Bernoulli de paramètre $p$.\\
    On pose $S_n = \sum_{i=1}^n X_i$ ($S_n \sim \text{B}(n,p)$).\\
    On a $Var(X_i) = p(1-p)$ pour tout $i$, et $Cov(X_i, X_j) = 0$ pour $i \neq j$.\\
    Donc : $Var(S_n) = \sum_{i=1}^n Var(X_i) = np(1-p)$.
}

\section{Loi jointe, loi marginale}

Soient $X_1, \ldots, X_d$ des variables aléatoires réelles à valeur dans des espaces dénombrables $E_1, \ldots, E_d$. Notons $(X_1, \ldots, X_d)$ la variable aléatoire à valeur dans $E_1 \times \cdots \times E_d$ définie par $(X_1, \ldots, X_d)(\omega) = (X_1(\omega), \ldots, X_d(\omega))$.\\

\definition{
    Soient $X_1, \ldots, X_d$ comme ci-dessus. On appelle \textbf{loi jointe} de $X_1, \ldots, X_d$ la loi de la variable aléatoire $(X_1, \ldots, X_d)$, c'est-à-dire la mesure de probabilité sur $E_1 \times \cdots \times E_d$ définie par :
    \[
        \mathbb{P}(X_1 = x_1, \ldots, X_d = x_d), \; \text{pour } x_1 \in E_1, \ldots, x_d \in E_d
    \]
}

\remark{La loi jointe de $X_1, \ldots, X_d$ contient toute l'information sur la distribution conjointe de ces variables.}

\definition{
    Soient $X_1, \ldots, X_d$ comme ci-dessus. On appelle \textbf{loi marginale} de $X_i$ la loi de la variable aléatoire $X_i$, c'est-à-dire la mesure de probabilité sur $E_i$ définie par :
    \[
        \mathbb{P}(X_i = x_i), \; \text{pour } x_i \in E_i
    \]
}

\theorem{Proposition}{Calcul de la loi marginale}{false}{
    Soient $X_1, \ldots, X_d$ comme ci-dessus. La loi marginale de $X_i$ se calcule à partir de la loi jointe par :
    \[
        \mathbb{P}(X_i = x_i) = \sum_{\substack{x_1 \in E_1 \\ \vdots \\ x_{i-1} \in E_{i-1} \\ x_{i+1} \in E_{i+1} \\ \vdots \\ x_d \in E_d}} \mathbb{P}(X_1 = x_1, \ldots, X_d = x_d)
    \]
}

\noindent{\textbf{Preuve :}\\
    On a que $\{ X_i = x_i \} = \bigcup_{x_1\in E_1} \cdots \bigcup_{x_{i-1} \in E_{i-1}} \bigcup_{x_{i+1} \in E_{i+1}} \cdots \bigcup_{x_d \in E_d} \{ X_1 = x_1, \ldots, X_d = x_d \}$, et que ces événements sont disjoints, donc l'union est disjointe.\\
    On conclut par la $\sigma$-additivité de la probabilité.
}

\example{
    (pour $d=2$) Soient $X,Y$ deux variables aléatoires à valeurs dans $E, F$.\\\\
    La loi marginale de $X$ est donnée par $\mathbb{P}(X=x) = \sum_{y \in F} \mathbb{P}(X=x, Y=y)$ pour $x \in E$.\\
    De même, la loi marginale de $Y$ est donnée par $\mathbb{P}(Y=y) = \sum_{x \in E} \mathbb{P}(X=x, Y=y)$ pour $y \in F$.
}

\section{Indépendance et loi conditionnelle}

\theorem{Proposition}{}{false}{
    Soient $X$ et $Y$ deux variables aléatoires à valeurs dans des ensembles dénombrables $E$ et $F$.\\
    Ces variables sont indépendantes si et seulement si il existe deux fonctions positives $f,g$ telles que pour tout $x \in E$ et $y \in F$ on ait :
    \[
        \mathbb{P}(X=x, Y=y) = f(x)g(y)
    \]
    Dans ce cas, on a $\mathbb{P}(X=x) = \frac{f(x)}{\sum_{x' \in E} f(x')}$ et $\mathbb{P}(Y=y) = \frac{g(y)}{\sum_{y' \in F} g(y')}$.
}
\remark{Cette proposition est valable pour plus de deux variables aléatoires.}

\definition{
    Soient $X$ et $Y$ deux variables aléatoires à valeurs dans des ensembles dénombrables $E$ et $F$.\\
    La \textbf{loi conditionnelle} de $Y$ sachant $X=x$ est la mesure de probabilité sur $F$ définie par :
    \[
        \mathbb{P}(Y=y | X=x) = \frac{\mathbb{P}(X=x, Y=y)}{\mathbb{P}(X=x)}, \; \text{pour } y \in F
    \]
}

\section{Tableaux de lois jointes}

Il est souvent pratique de représenter la loi jointe de deux variables à valeurs dans des espaces finis par un tableau à double entrée.\\\\
Pour obtenir les lois marginales à partir de ce tableau, il suffit de faire la somme des probabilités sur les lignes (pour la loi marginale de $X$) ou sur les colonnes (pour la loi marginale de $Y$).\\\\
Pour obtenir la loi conditionnelle de $Y$ sachant $X=x$, il suffit de diviser les probabilités de la ligne correspondant à $X=x$ par la probabilité marginale $\mathbb{P}(X=x)$.

\newpage
\tableofcontents

\end{document}